\section{Evolución de los métodos de generación 3D}

\subsection{Métodos tempranos: 3D GAN (2016)}

Los primeros métodos de modelado 3D generativo utilizaron redes generativas adversariales (GAN). En el trabajo de 2016, los investigadores construyeron un modelo puramente generativo que partía de variables latentes ruidosas para entrenar un modelo a generar representaciones 3D en forma de vóxeles (voxel), utilizando un \textbf{discriminador 3D} para el entrenamiento adversarial.

Las dos limitaciones principales de este método son:
\begin{itemize}
    \item \textbf{Imposibilidad de controlar la entrada con precisión}: El proceso de generación es completamente aleatorio, por lo que no se puede especificar qué objeto concreto generar.
    \item \textbf{Necesidad de datos verdaderos en 3D}: El entrenamiento depende de ground truth en 3D, y los datos de alta calidad en 3D son extremadamente escasos.
\end{itemize}

\subsection{PlatonicGAN: Independencia de los datos 3D}

La innovación central de PlatonicGAN radica en la introducción de un \textbf{renderizador diferenciable} (differentiable renderer) y un \textbf{discriminador 2D}, eliminando así la dependencia de los datos verdaderos en 3D. Esto permite entrenar el modelo directamente sobre datos de imágenes 2D disponibles en Internet. Además, PlatonicGAN introduce un \textbf{codificador de imágenes}, logrando la capacidad de reconstrucción 3D a partir de una sola imagen.

\begin{knowledgebox}{Significado del renderizado diferenciable}
El renderizado diferenciable (differentiable rendering) actúa como puente entre las representaciones 3D y las señales de supervisión en 2D. Al hacer el proceso de renderizado diferenciable, los gradientes se propagan hacia atrás desde la función de pérdida en el espacio de imágenes 2D hasta los parámetros de la representación 3D, permitiendo entrenar modelos generativos 3D con grandes volúmenes de datos de imágenes 2D sin necesidad de costosas anotaciones en 3D.
\end{knowledgebox}

\subsection{PixelNeRF: Fusión multivista}

Sobre la base del enfoque de reconstrucción basado en imágenes, PixelNeRF se centra en la \textbf{entrada multivista}. Su innovación central consiste en agregar (aggregate) las características codificadas provenientes de múltiples imágenes de entrada y fusionarlas en una representación 3D coherente: específicamente, un NeRF (Neural Radiance Field).

\subsection{Características comunes y limitaciones}

Todos los métodos anteriores comparten un rasgo común: predicen directamente una representación \textbf{explícita y consistente en 3D}.

\begin{importantbox}{Ventajas y desventajas de la representación explícita 3D}
\textbf{Ventaja}: Una vez obtenido el modelo 3D, se puede renderizar rápidamente y a bajo costo desde cualquier punto de vista, lo que es ideal para aplicaciones en tiempo real en dispositivos móviles como teléfonos inteligentes y tabletas.

\textbf{Desventaja}: La calidad general de generación es baja. Existen dos razones: (1) los datos de alta calidad en 3D son extremadamente escasos; (2) estos métodos no deberían utilizar un sesgo inductivo demasiado fuerte (como el renderizado), pero evitar dicho sesgo limita las capacidades del modelo.
\end{importantbox}

\subsection{Resumen de la sección}

Los métodos de generación 3D evolucionaron desde los 3D GAN, que dependían de datos verdaderos en 3D, hasta PlatonicGAN, que aprovechaba el renderizado diferenciable y el discriminador 2D, y posteriormente a PixelNeRF, que admite entradas multivista. La ventaja común de estos métodos radica en la generación de representaciones explícitas en 3D y un bajo costo de renderizado; sin embargo, debido a la escasez de datos en 3D, su calidad general es inferior a la de los modelos generativos en 2D.
